\documentclass[11pt]{article}
\usepackage{graphicx}
\setlength{\parindent}{0pt}
\usepackage{hyperref}
\usepackage{enumitem}
\usepackage[utf8]{inputenc}
\usepackage[T1]{fontenc}
\usepackage[spanish]{babel}
\usepackage[left=1.06cm,top=1.7cm,right=1.06cm,bottom=1.0cm]{geometry}
\hypersetup{colorlinks=true, urlcolor=black, linkcolor=black}
\pagestyle{empty}

% --- Formato de subtítulos de sección (sin numeración, centrados, con línea) ---
\newcommand{\seccion}[1]{%
    \vspace{4pt}
    \begin{center}
        \textbf{\MakeUppercase{#1}}
    \end{center}
    \vspace{-10pt}
    \hrulefill
    \vspace{2pt}
}
\setlength{\parindent}{10pt}
\begin{document}

% ============ ENCABEZADO ============
\begin{center}
    {\LARGE \textbf{Ernesto Mendieta Cuecuecha}}\\[2pt]
    Ingeniero en Inteligencia Artificial\\[4pt]
    Santa Ana Chiautempan, Tlaxcala, México \textbullet\ 222 825 1582 \textbullet\ \href{mailto:ernestomecue@gmail.com}{ernestomecue@gmail.com}\\
    \href{https://www.linkedin.com/in/ernesto-mendieta-cuecuecha-90135125a/}{linkedin.com/in/ernesto-mendieta-cuecuecha} \textbullet\ \href{https://github.com/ErnestoMendieta}{github.com/ErnestoMendieta}
\end{center}

% ============ PERFIL PROFESIONAL ============
\seccion{Perfil Profesional}

\hspace{\parindent} Ingeniero en Inteligencia Artificial egresado del Instituto Politécnico Nacional, con experiencia aplicada en aprendizaje automático, LLM's y visión artificial. Dominio de Python, ElevenLabs y Hugging Face, con experiencia en desarrollo de sistemas de llamadas multiagente de inteligencia artifical y comportamiento de LLM's. Historial comprobado en proyectos de detección de objetos, análisis de sentimientos y segmentación semántica, con enfoque en resultados medibles y trabajo colaborativo bajo metodologías ágiles. Uso de Agentes de código como Codex y Claude Code.

% ============ EDUCACIÓN ============
\seccion{Educación}

\textbf{Instituto Politécnico Nacional (UPIIT)} \hfill Tlaxcala, México\\
Ingeniería en Inteligencia Artificial \hfill Agosto 2021 -- Julio 2025

\vspace{8pt}

\textbf{Centro de Bachillerato Tecnológico Industrial y de Servicios No. 03} \hfill Tlaxcala, México\\
Técnico en Programación -- Cédula Profesional 13982196 \hfill Agosto 2019 -- Julio 2021

% ============ EXPERIENCIA LABORAL ============
\seccion{Experiencia Laboral}

\textbf{Secretariado Ejecutivo del Sistema Estatal de Seguridad Pública (SESESP)} \hfill Tlaxcala de Xicotencatl\\
\textbf{Analista de Inteligencia Artificial} \hfill Enero 2026 -- Actualidad
\begin{itemize}[noitemsep, topsep=2pt, partopsep=0pt, parsep=2pt, leftmargin=0.55cm]
    \item Implementé herramientas de inteligencia artificial en la extracción de identidades en narrativas de incidentes.
    \item Desarrollo de sistema multiagente de voz para la atención a la ciudadanía.
    \item Uso de mejores prácticas de ingeniería de prompts e ingeniería de contexto en agentes de Inteligencia Artificial.
\end{itemize}

\vspace{8pt}

\textbf{Centro de Innovación y Desarrollo Tecnológico en Cómputo (CIDETEC -- IPN)} \hfill Ciudad de México, México\\
\textbf{Residente de Prácticas Profesionales -- Visión Artificial} \hfill Enero 2025 -- Mayo 2025
\begin{itemize}[noitemsep, topsep=2pt, partopsep=0pt, parsep=2pt, leftmargin=0.55cm]
    \item Desarrolló un proyecto de visión artificial para la segmentación semántica de objetos en entornos urbanos.
    \item Ajusté (\textit{fine-tuning}) un modelo Transformer con mecanismo de atención sobre el conjunto de datos CityScapes, mejorando la precisión en la identificación de objetos.
    \item Implementó el proyecto en Python utilizando cómputo paralelo y el framework Hugging Face para optimizar el rendimiento del modelo.
\end{itemize}

\vspace{8pt}

\textbf{CETis 132} \hfill Tlaxcala, México\\
\textbf{Profesor de Cálculo Diferencial e Integral} \hfill Septiembre 2024 -- Diciembre 2024
\begin{itemize}[noitemsep, topsep=2pt, partopsep=0pt, parsep=2pt, leftmargin=0.55cm]
    \item Impartió clases de Cálculo Diferencial e Integral a estudiantes de nivel medio superior, fortaleciendo su comprensión de los fundamentos de la materia.
    \item Diseñó material didáctico apoyado en inteligencia artificial generativa para facilitar la comprensión de conceptos complejos.
    \item Automatizó tareas administrativas y de gestión de clases mediante scripts en Python e integración con Google Classroom, optimizando el tiempo de preparación y calificación.
\end{itemize}
 \newpage
% ============ PROYECTOS ============
\seccion{Proyectos}

\textbf{Detección de Riesgos Urbanos para Personas con Discapacidad Visual}
\begin{itemize}[noitemsep, topsep=2pt, partopsep=0pt, parsep=2pt, leftmargin=0.55cm]
    \item Entrenó un modelo de visión artificial para identificar situaciones de riesgo en entornos urbanos, priorizando bajos tiempos de inferencia y respuesta para asistencia a personas con discapacidad visual.
\end{itemize}

\vspace{6pt}

\textbf{Chatbot Musical} -- \href{https://github.com/ErnestoMendieta/chatbot-musical-con-lstm}{github.com/ErnestoMendieta/chatbot-musical-con-lstm} \hfill
\begin{itemize}[noitemsep, topsep=2pt, partopsep=0pt, parsep=2pt, leftmargin=0.55cm]
    \item Diseñó un chatbot de recomendación musical basado en redes LSTM para modelado de lenguaje y clasificación de sentimientos.
    \item Integró la API de Spotify mediante autenticación OAuth y una aplicación móvil desarrollada en Flutter como interfaz de usuario.
\end{itemize}

\vspace{6pt}

\textbf{GunGuardAI} -- \href{https://github.com/ErnestoMCUpiit/DeteccionArmasDeFuego/tree/foo}{github.com/ErnestoMCUpiit/DeteccionArmasDeFuego} \hfill
\begin{itemize}[noitemsep, topsep=2pt, partopsep=0pt, parsep=2pt, leftmargin=0.55cm]
    \item Desarrolló una aplicación móvil en Flutter para la detección de armas de fuego mediante un modelo SSD (Single Shot Detection) preentrenado.
    \item Implementó detección en tiempo real a partir de fotografía y de video capturado con la cámara del dispositivo.
\end{itemize}

\vspace{6pt}

\textbf{Flotilla App} -- \href{https://github.com/ErnestoMCUpiit/flotillaWeb}{github.com/ErnestoMCUpiit/flotillaWeb} 
\hfill
\begin{itemize}[noitemsep, topsep=2pt, partopsep=0pt, parsep=2pt, leftmargin=0.55cm]
    \item Construyó una aplicación web y móvil (Flutter) para la gestión de flotillas de vehículos de transporte, con acceso y sincronización de datos en tiempo real mediante Firebase.
\end{itemize}

\vspace{6pt}

\textbf{Seguidor de Trayectoria con Visión Artificial}
\begin{itemize}[noitemsep, topsep=2pt, partopsep=0pt, parsep=2pt, leftmargin=0.55cm]
    \item Construyó un seguidor de línea con microcontroladores Arduino y ESP32, utilizando un modelo YOLO preentrenado para clasificar señales de dirección ("izquierda", "derecha", "alto") capturadas por cámara y generar instrucciones de navegación en tiempo real.
\end{itemize}

% ============ HABILIDADES TÉCNICAS ============
\seccion{Habilidades Técnicas e Idiomas}

\textbf{Lenguajes de programación:} Python (intermedio), JavaScript (básico)\\
\vspace{6pt}
\textbf{Frameworks y herramientas:} TensorFlow, Hugging Face, Flutter, React / React Native, Firebase, Pandas, ElevenLabs, Git / GitHub\\
\vspace{6pt}
\textbf{Idiomas:} Español (nativo); Inglés -- Nivel B1, certificado por el Centro de Lenguas Extranjeras del IPN (CENLEX)

\end{document}